% !Mode:: "TeX:UTF-8"
% !TEX program  = xelatex
%
% SlimeVR DIY 全身动作捕捉项目 —— 单一源文档
% 所有后续设计、实现与评测记录均追加于此文件。
\documentclass[UTF8,11pt,a4paper]{ctexart}

\usepackage[margin=1in]{geometry}
\usepackage{amsmath, amssymb, amsthm}
\usepackage{booktabs}
\usepackage{multirow}
\usepackage{makecell}
\usepackage{graphicx}
\usepackage{float}

% 架构图（单色/中性，不使用彩色）
\usepackage{tikz}
\usetikzlibrary{shapes.geometric, arrows.meta, positioning, calc, fit, backgrounds}
\tikzstyle{anchorbox} = [rectangle, draw=black, thick, rounded corners, fill=gray!12, minimum height=1.0cm, align=center, font=\small]
\tikzstyle{solverbox} = [rectangle, draw=black, thick, rounded corners, fill=white, minimum height=1.0cm, align=center, font=\small]
\tikzstyle{imubox}    = [rectangle, draw=black, thick, rounded corners, fill=gray!6,  minimum height=0.9cm, align=center, font=\small]
\tikzstyle{residbox}  = [rectangle, draw=black, thick, rounded corners, fill=white, minimum height=0.9cm, align=center, font=\small]
\tikzstyle{arrow}     = [-{Stealth}, thick, draw=black]

% 算法环境（中文支持）
\usepackage{algorithm}
\usepackage{algorithmicx}
\usepackage{algpseudocode}
\floatname{algorithm}{算法}
\renewcommand{\algorithmicrequire}{\textbf{输入:}}
\renewcommand{\algorithmicensure}{\textbf{输出:}}

% 代码清单（单色）
\usepackage{listings}
\lstset{
    basicstyle=\small\ttfamily,
    frame=single,
    breaklines=true,
    numbers=left,
    numberstyle=\tiny,
    showstringspaces=false,
    keywordstyle=\bfseries,
    commentstyle=\itshape,
    language=Python
}

% 链接（隐藏颜色，保持单色）
\usepackage[hidelinks]{hyperref}

\title{\LARGE\bfseries 基于 SlimeVR 的低成本全身动作捕捉\\与服务器端神经网络误差纠正}

\author{%
    小组组长：\underline{\hspace{2.5cm}} \\
    小组成员：姓名+学号
}

\date{\today}

\begin{document}

\maketitle

\begin{abstract}
本项目基于 SlimeVR 开源方案，DIY 一套以 nRF52840 + LSM6DSV16X 为核心、覆盖胸/髋/两大腿/两小腿的 6 点全身动作捕捉系统，用于 VRChat 虚拟形象驱动。核心创新是在 PC 服务端引入一层轻量神经网络，作为“姿态先验 / 去噪器”，对多传感器姿态做联合误差纠正，弥补官方服务端仅依赖骨骼几何约束与最小二乘自动测比例的不足。针对 6 轴 IMU 无磁力计、偏航（yaw）随时间累积漂移的问题，本设计采用“上半身绝对锚 + 下半身残差兜底”的分层抑制架构：上半身由 Virtual Desktop Estimation 虚拟追踪器（绝对偏航、零漂移）锚定，下半身局部小漂由 Stay Aligned（规则式）、本设计神经网络残差校正器与骨架几何约束三者共同吸收。训练数据完全合成——以 AMASS 等公开动捕数据为干净监督，注入 LSM6DSV16X 量级的偏航随机游走漂移构造 corrupted→clean 配对，无需真人标注。端到端管线延迟目标 < 20\,ms，硬件总成本约 ¥655。
\end{abstract}

\noindent\textbf{文档说明：}本文件为 SlimeVR DIY 全身动捕项目的单一源文档，后续所有设计修订、实现记录与评测结果均追加于此，不另建独立文档。

\newpage
\tableofcontents
\newpage

% ==================== 一、项目内容 ====================
\section{一、项目内容}
本项目基于 SlimeVR 开源方案，DIY 一套以 nRF52840 + LSM6DSV16X 为核心的 6 点全身动作捕捉系统，用于 VRChat 虚拟形象驱动；核心创新是在 PC 服务端引入轻量神经网络，对多传感器姿态做联合误差纠正，弥补官方仅靠几何约束纠偏的不足。

% ==================== 二、可行性 ====================
\section{二、项目可行性分析}
\textbf{必要性：}商用全身追踪方案价格高昂，而 SlimeVR 开源硬件可将整套成本压到千元以内，降低虚拟形象驱动门槛。

\textbf{创新性：}官方服务端仅依赖骨骼几何约束与最小二乘自动测比例（AutoBone）与脚锁做纠偏，本设计加入神经网络姿态先验，可抑制抖动、穿模与偏航漂移。

\textbf{人员：}团队具备嵌入式开发与机器学习基础，分工明确、可覆盖硬件到算法。

\textbf{时间：}采购、焊接、固件编译、模型训练与联调均可在学期内完成。

% ==================== 三、问题 ====================
\section{三、可能需要解决的问题}
\textbf{硬件与固件：}nRF 固件对 LSM6DSV16X 的识别需实机验证（已设闸门 G1），且 nRF Connect SDK 编译环境庞大、追踪器与接收器需使用不同版本，搭建易出错。

\textbf{工程：}焊接与烧录顺序错误可能导致设备变砖，已准备 DAPLink 救砖。

\textbf{算法：}神经网络缺乏带噪动捕训练数据，拟用 AMASS 等公开数据集合成带噪监督信号；服务端为 Java，模型推理需经 OSC/VMC 松耦合桥接或 fork 接入，须将端到端延迟控制在 20\,ms 以内。

% ==================== 四、材料 ====================
\section{四、需要购买的材料}

\begin{table}[H]
\centering
\small
\caption{项目材料清单与预算（合计 ¥655）}
\label{tab:materials}
\begin{tabular}{clr}
\toprule
\textbf{序号} & \textbf{材料名} & \textbf{价格} \\
\midrule
1 & nRF52840 主控板 + LSM6DSV16X IMU（6 点及备件） & ¥370 \\
2 & 电池 / 无线接收器 / DAPLink / 外壳 / 工具 & ¥285 \\
\midrule
\multicolumn{2}{c}{\textbf{合计}} & \textbf{¥655} \\
\bottomrule
\end{tabular}
\end{table}

% ==================== 五、实施计划 ====================
\section{五、项目实施计划}

\subsection{5.1 人员任务分配}
组长：负责总体方案制定、nRF 固件编译与烧录、SlimeVR 服务端集成与联调。成员 A：负责服务器端神经网络模型的设计、训练与 OSC/VMC 桥接实现。成员 B：负责硬件采购、焊接组装、3D 打印外壳与项目文档撰写。

\subsection{5.2 计划时间表}
第 1–2 周：搭建开发环境，跑通 SlimeVR 软件四链（Server / GUI / OpenVR 驱动）。第 3 周：采购物料并完成 6 点 tracker 焊接组装。第 4 周：烧录固件，实机完成 G1 验证（IMU 识别）。第 5–6 周：设计并训练服务器端神经网络误差纠正模型。第 7–8 周：模型与 SlimeVR 集成，VRChat 实机测试与调参。第 9 周：整理文档、准备答辩。

% ==================== 六、NN 设计 ====================
\section{六、服务器端神经网络模型设计}

\subsection{6.1 问题定位}
SlimeVR 官方服务端仅用骨骼几何约束、最小二乘自动测比例（AutoBone）与脚锁做纠偏，各 tracker 的绝对朝向由固件侧 VQF 融合给出，服务端信任该朝向并装配成骨架。其局限是：无跨 tracker 的统计融合、不自动纠正偏航（yaw）漂移、单点失稳仅由相邻关节约束兜底。本设计在 PC 服务端增加一层轻量神经网络，作为“姿态先验 / 去噪器”，对多传感器姿态做联合误差纠正。

\subsection{6.2 输入与输出}
\textbf{输入：}各 tracker 的绝对朝向四元数 $q$ 与角速度 $\omega$（来自固件 VQF 输出），以及用户骨长向量；采用长度为 $N$ 的滑动窗口（近 $N$ 帧）以利用时序一致性。

\textbf{输出：}纠正后的全身关节位置 / 朝向，或直接输出各 tracker 朝向的残差校正量；下游仍以骨长约束与 HMD 锚定作最终装配。

\subsection{6.3 网络结构}
采用轻量时序模型：双向 LSTM / TCN，或小型 Transformer 编码器；亦可建模为 VAE 式姿态先验，用 KL 正则把输出拉向合理人体姿态流形。参数量控制在百万级以内，推理用 FP16，确保可在消费级 GPU 实时运行。

\subsection{6.4 训练数据与目标函数}
\textbf{数据：}以 AMASS 等公开动捕数据集为干净监督，按 SlimeVR 噪声特性（陀螺积分漂移、yaw 缓漂、偶发掉点）合成带噪骨架作为网络输入。

\textbf{损失：}位置 / 朝向回归误差 + 物理合理性项（骨长守恒、关节角限、地面约束）；亦可用无监督对抗 / 重构学习人体姿态流形。

\subsection{6.5 与 SlimeVR 的集成方式}
\textbf{松耦合（推荐）：}SlimeVR 经 OSC / VMC 输出骨架 → 独立 NN 进程消费 → 清洗后回写 OSC / VMC 给 VRChat，不改动 SlimeVR 代码。

\textbf{紧耦合：}fork SlimeVR Server，在骨架计算后插入推理；需经 JNI / 子进程 / ONNX Runtime 桥接 Java 与模型。

\subsection{6.6 延迟与部署约束}
VR 要求端到端管线 < 20\,ms；NN 推理须 < 5\,ms（小模型 + TensorRT / ONNX Runtime）。PC 端 GPU 完全可行；固件侧 nRF52840 无 NPU、算力不足，模型不应下沉到设备端。

\subsection{6.7 验证方案}
在真实 tracker 采集数据与合成带噪数据上，对比加入 NN 前后的抖动均方根、穿模率、偏航漂移量与用户主观评分，验证纠正效果与实时性。

\subsection{6.8 全系统集成架构（VD Estimation 虚拟追踪器锚定上半身 + 6 轴 IMU 下半身）}
本项目 6 个 tracker 分布于胸、髋、两大腿、两小腿，全部采用 6 轴 LSM6DSV16X（加速度计 + 陀螺仪，无磁力计）。单靠 IMU 无法获得绝对偏航参考：加速度计只能靠重力校正俯仰与横滚，偏航（绕重力轴的旋转）只能由陀螺仪积分得到，其零偏与角随机游走随时间累积（约 $0.1^{\circ}$/s 量级即 $6^{\circ}$/min），表现为全身缓慢转偏。为此本设计引入“上半身绝对锚 + 下半身残差兜底”的分层抑制架构。

\subsubsection{6.8.1 组件与数据流}
上半身姿态（胸 / 肩 / 肘 / 腕）由 Virtual Desktop 的 Estimation 功能提供：它利用 Quest 3 头显自带的下视侧摄做 Inside-Out Body Tracking（IOBT），生成一组“虚拟 tracker”，其位姿为相机绝对测量、偏航零漂移，且零额外硬件。VD 虚拟 tracker 在 SteamVR 中以设备形式暴露，经由官方 SlimeVR Feeder App（随 SlimeVR 安装器附带、随 SteamVR 启动的 OpenVR 桥）通过命名管道喂入 SlimeVR-Server，与本项目的 6 个 IMU 在同一套骨架求解器中统一解算。社区工具 Virtual Desktop Body Tracking Configurator 可关闭 VD 的下半身虚拟 tracker（Generative Legs），仅保留上半身，把腿留给 IMU。

完整数据流为：IMU trackers 经 Wi-Fi 入 SlimeVR-Server 做骨架解算；VD Estimation 虚拟 tracker 经 SteamVR 设备由 Feeder App 桥接进同一 Server；两者在同一套骨架求解器中统一解算后，经 OSC/VMC 输出至 VRChat。官方兼容文档明确 SlimeVR 与 Virtual Desktop Estimation “specifically for arms” 搭配使用。

\subsubsection{6.8.2 漂移抑制分层}
第一层（全局锚，消除全局偏航漂移）：VD Estimation 的胸 / 上脊柱（可含髋）为绝对参考，骨架求解器自胸向下做前向运动学传播，胸一经固定，原本会贯穿全身的全局偏航漂移即被砍除。第二层（局部残差，吸收下半身小漂）：仅余大腿 / 小腿的局部偏航小漂，幅度远小于全局漂，由以下三者吸收——(a) SlimeVR 内置 Stay Aligned（基于校准放松姿态的姿态先验，规则式梯度下降）；(b) 本设计 6.1–6.7 的 NN 残差校正器；(c) 骨架几何约束（骨长守恒、关节角限、地面约束、相邻 tracker 互推）。

\begin{figure}[H]
\centering
\begin{tikzpicture}[
    node distance=1.6cm and 3.2cm
]
\node[anchorbox] (vd) {VD Estimation 虚拟追踪器\\（胸/上脊柱，绝对偏航，零漂移）};
\node[solverbox] (solver) [below=of vd] {SlimeVR 骨架求解器\\（统一解算 IMU + VD 虚拟追踪器）};
\node[imubox]    (imu)    [left=of solver] {下半身 IMU trackers\\（大腿/小腿，局部偏航小漂）};
\node[residbox]  (resid)  [right=of solver] {残差吸收层（每 tick）：\\Stay Aligned（规则式）\\NN 残差校正器（本设计）\\骨架几何约束};
\node[solverbox] (vrc)    [below=of solver] {OSC / VMC 输出 $\to$ VRChat};

\draw[arrow] (vd) -- (solver) node[midway, left] {固定全局偏航};
\draw[arrow] (imu) -- (solver) node[midway, above] {局部小漂};
\draw[arrow] (solver) -- (vrc) node[midway, left] {骨架};
\draw[<->, arrow] (resid) -- (solver) node[midway, above] {每 tick 残差吸收};
\end{tikzpicture}
\caption{分层漂移抑制架构：上半身绝对锚（VD）消除全局偏航漂移，下半身局部小漂由残差吸收层兜底}
\label{fig:layered_arch}
\end{figure}

\subsubsection{6.8.3 Stay Aligned 与 NN 的互补}
Stay Aligned 是规则式姿态先验，每 tick 将 tracker 偏航朝“校准时的放松姿态”梯度下降微调，不读取头显偏航，本身不是绝对参考；其盲区在于非放松姿态（侧躺、大幅跳舞、前倾）时放松姿态不成立而失效。NN 残差校正器从 AMASS 等大规模动捕学到“人类在各活动下的合理姿态分布”，可覆盖上述盲区。二者同属姿态先验类，NN 是 Stay Aligned 的补强而非替代；代价是 NN 可能“幻觉”（将姿态掰到合理但错误的位置），故必须做残差 + 置信门控（见 6.9）。

\subsection{6.9 残差校正器的实现纪律与评测}
为使 NN 稳妥，须遵守以下纪律：(1) 仅输出小幅度残差四元数 $\Delta q$ 并叠加回原朝向，绝不整体替换骨架；(2) 加置信门控，低置信时不修正、退回原值；(3) 训练数据全合成——以 AMASS 干净全身骨架为监督，经简化前向运动学重建后，在下半身关节注入 LSM6DSV16X 量级的偏航随机游走漂移（陀螺零偏 / 角随机游走）与 IMU 噪声，构造 corrupted→clean 配对回归残差，无需真人标注；(4) 结构以 TCN / 双向 LSTM 起步，Transformer 编码器作后续升级；(5) 评测先量化后上线——在合成漂移上测均方误差与偏航误差恢复率（如注入 $10^{\circ}$/min 漂，看 NN 收敛回多少），再上真机对比 reset 频率与穿模率，不依赖“感觉更稳”的主观判断。

\noindent\textbf{算法 1（推理，每帧）：}轻量残差校正器在骨架求解后的作用流程。

\begin{algorithm}[H]
\caption{NN 残差校正器推理流程（每帧）}
\label{alg:infer}
\begin{algorithmic}[1]
\REQUIRE 滑动窗口 $\mathcal{W}=\{(q_t,\omega_t)\}_{t=1}^{N}$（各 tracker 局部朝向与角速度）；VD 胸/上脊柱绝对锚姿 $T_{\text{VD}}$；骨架约束 $\mathcal{C}$
\ENSURE 纠正后下半身关节朝向 $q^{\text{corr}}$
\STATE 若 $T_{\text{VD}}$ 可用：取全局偏航参考 $\psi_{\text{ref}} = \text{yaw}(T_{\text{VD}})$，在求解器中锁定全身全局偏航
\STATE 读取 IMU 下半身 tracker 当前局部朝向 $q$ 与角速度 $\omega$，压入窗口 $\mathcal{W}$
\STATE 前向 NN（TCN / 双向 LSTM）：$(\Delta q, c) = f_{\theta}(\mathcal{W})$，其中 $c\in[0,1]$ 为置信度
\IF{$c < \tau$}
    \STATE $\Delta q \leftarrow \text{identity}$（低置信：不修正，退回原值）
\ENDIF
\STATE 叠加残差：$q^{\text{corr}} = \Delta q \otimes q$（仅小幅修正，不整体替换）
\STATE 施加骨架约束 $\mathcal{C}$（骨长守恒、关节角限、地面约束、相邻 tracker 互推）
\STATE 经 OSC / VMC 输出 $q^{\text{corr}}$ 至 VRChat
\RETURN $q^{\text{corr}}$
\end{algorithmic}
\end{algorithm}

\noindent\textbf{算法 2（训练，离线）：}全合成 corrupted→clean 配对构造与回归。

\begin{algorithm}[H]
\caption{NN 残差校正器训练流程（离线、全合成）}
\label{alg:train}
\begin{algorithmic}[1]
\REQUIRE AMASS 干净动捕序列；LSM6DSV16X 噪声配置（陀螺零偏 $b$、角随机游走 $N_{\text{ARW}}$）；窗口长度 $N$
\ENSURE 模型参数 $\theta$
\STATE 由 AMASS 干净骨架经简化前向运动学重建干净下半身局部朝向 $q^{*}_{\text{clean}}$
\STATE 注入合成漂移：对每个下半身关节，以随机游走叠加陀螺零偏生成偏航漂移，并叠加 IMU 噪声，得到 corrupted 窗口 $\mathcal{W}^{\text{corrupt}}$
\STATE 计算监督残差目标 $\Delta q^{*} = q^{*}_{\text{clean}} \otimes q^{-1}_{\text{corrupt}}$
\STATE 构造数据集 $\{(\mathcal{W}^{\text{corrupt}}_i, \Delta q^{*}_i)\}_{i}$
\STATE 训练 TCN / 双向 LSTM：$(\Delta q, c) = f_{\theta}(\mathcal{W})$；损失
\[
\mathcal{L} = \text{SmoothL1}(\Delta q, \Delta q^{*}) + \lambda_{\text{phys}}\,\mathcal{L}_{\text{phys}} + \lambda_{c}\,\mathcal{L}_{\text{conf}}
\]
其中 $\mathcal{L}_{\text{phys}}$ 为骨长守恒 / 关节角限正则，$\mathcal{L}_{\text{conf}}$ 为置信校准项
\STATE 在合成漂移上验证：注入 $10^{\circ}$/min 漂，测偏航误差恢复率；仅当达标后再上真机
\RETURN $\theta$
\end{algorithmic}
\end{algorithm}

\noindent\textbf{实现骨架（PyTorch 伪代码，单色清单）：}

\begin{lstlisting}[caption={TCN 残差校正器与合成漂移注入（实现骨架）}, label={lst:skeleton}]
import torch, torch.nn as nn

class TCNResidualCorrector(nn.Module):
    # 输入: (B, N, J*8)  窗口帧 x (关节 local q(4)+omega(3)+bone(1))
    # 输出: (B, J, 4) 残差四元数 + (B, J) 置信度
    def __init__(self, in_ch, hidden, J, levels=4):
        super().__init__()
        layers = []
        for i in range(levels):
            dil = 2 ** i
            layers += [nn.Conv1d(in_ch, hidden, 3, dilation=dil, padding=dil),
                       nn.ReLU(), nn.Dropout(0.1)]
        self.tcn = nn.Sequential(*layers)
        self.head_q = nn.Conv1d(hidden, J * 4, 1)
        self.head_c = nn.Conv1d(hidden, J, 1)

    def forward(self, x):
        h = self.tcn(x.transpose(1, 2))          # (B, hidden, N)
        dq = self.head_q(h).view(-1, J, 4)        # 残差四元数(未归一)
        dq = dq / dq.norm(dim=-1, keepdim=True)
        c = torch.sigmoid(self.head_c(h)).view(-1, J)
        return dq, c

def inject_synthetic_drift(q_clean, b, narw, dt):
    # 随机游走偏航漂移 + IMU 噪声，模拟 LSM6DSV16X 量级
    psi = torch.cumsum(b * dt + narw * torch.randn_like(q_clean[..., 0]) * dt ** 0.5, 0)
    q_drift = yaw_ quaternion(psi)               # 仅绕重力轴旋转
    q_corrupt = q_drift * q_clean
    return q_corrupt

# 训练步（节选）
dq_pred, conf = model(window_corrupt)
loss = smooth_l1(dq_pred, dq_target) \
     + lam_phys * physics_reg(dq_pred, bone_len) \
     + lam_conf * confidence_calib(conf, dq_err)
loss.backward(); opt.step()
\end{lstlisting}

\subsection{6.10 现有方案对比与项目定位}
经社区与学术检索，SlimeVR 生态内\textbf{不存在}专门的“神经网络残差校正器”开源项目，这一方向属生态空白；官方现有漂移缓解手段为 Stay Aligned（规则式）+ 偏航重置（yaw reset）+ 可选磁力计。学术界有面向单 IMU 的 AHRS 神经网络基线（RoNIN、TLIO），但其未与 SlimeVR 集成；视觉 ML 全身追踪（KinectAIFusionFBT、Fever、Mediapipe-VR、ToucanTrack）则提供另一条绝对参考路线。本设计定位如下表。

\begin{table}[H]
\centering
\small
\caption{现有方案与本设计的对比定位}
\label{tab:compare}
\begin{tabular}{p{3.0cm}p{2.6cm}p{2.0cm}p{4.0cm}}
\toprule
\textbf{方案} & \textbf{类型} & \textbf{是否绝对参考} & \textbf{局限 / 本项目定位} \\
\midrule
Stay Aligned & 规则式姿态先验 & 否（相对校准放松姿态） & 非放松姿态（侧躺/大幅舞蹈）失效；\textbf{保留为基线} \\
RoNIN / TLIO & 学术 AHRS 神经网络（单 IMU） & 否 & 单 IMU、无全身约束、需真机标定；\textbf{借鉴其网络思想} \\
相机 ML 全身追踪（KinectAIFusionFBT / Fever / Mediapipe-VR / ToucanTrack） & 视觉绝对参考 & 是 & 需额外摄像头/算力、遮挡敏感；\textbf{替代路线，本设计不采用} \\
\textbf{本设计}（VD 锚 + NN 残差校正） & 分层（绝对锚 + 统计残差） & 上半身绝对（VD）、下半身残差 & NN 幻觉需置信门控；\textbf{主路线} \\
\bottomrule
\end{tabular}
\end{table}

\end{document}
